% 3.7: abstract, at most 500 words.
\FrontChapter{Abstract}
\thispagestyle{empty}

Swarms of unmanned aerial vehicles (UAVs) can learn a shared perception model with
federated learning (FL): each UAV trains on the images it collects and sends only a
model update to a ground base station. In a tactical setting two problems arise
together. Air-to-ground links lose packets because of fading, distance, handover and
jamming, and some UAVs may be captured and made to send poisoned updates. Most
Byzantine-robust aggregation rules assume that every update arrives intact, while
most work on FL over unreliable links assumes that every client is honest. This
project studies both problems in one system.

We built a simulation framework in PyTorch for hierarchical FL over a UAV swarm. Ten
UAVs with non-identically distributed aerial images are served by three base
stations, which aggregate locally and synchronize every five rounds. Uplink packet
losses come either from a synthetic channel with path loss, Rayleigh fading and a
ground jammer, or from a measured UAV received-power trace (DeepSense~6G Scenario~23)
replayed at a chosen signal-to-noise ratio. Three of the UAVs mount sign-flip,
ALIE-style or inner-product-manipulation attacks. The framework includes FedAvg,
coordinate-wise median, FLTrust, a cosine-weighted rule taken from prior UAV-FL
work, and a reputation-based aggregator, CCRA, that we designed. CCRA compares each
update with a reference update computed on a small trusted dataset, keeps a
reputation for every UAV across rounds, and can use link information in two ways:
by ignoring coordinates whose packets were lost (erasure masking), and by letting
link telemetry slow down how fast a UAV regains trust.

The work proceeded in three phases. A FedAvg baseline reached
62.6\% (IID) and 60.2\% (non-IID) test accuracy. A preliminary single-seed study
suggested that CCRA resisted a label-flipping attack better than the cosine-weighted
baseline. To test which part of CCRA was responsible, we then ran a controlled,
prespecified study of \NumRuns{} training runs on an A100 GPU: \NumSeeds{} paired
seeds, four attack settings, three channel conditions and six aggregators. Diverged
runs were scored at chance rather than discarded, and all uncertainty was estimated
with a paired bootstrap over seeds.

The study gave a clear and partly negative answer. On a weak (5\,dB) link under
attack, erasure masking improved accuracy by \MaskGain{} percentage points (95\%
confidence interval \ci{\MaskLo}{\MaskHi}). Adding link telemetry to the reputation
update changed accuracy by \TelGain{} points (\ci{\TelLo}{\TelHi}) because it
rejected \RejCcra\% of honest updates, against \RejMask\% without telemetry. We
explain this with a short analysis of the reputation rule: when rewards are slowed
but penalties are not, an honest UAV's reputation settles near its worst rounds. All
reference-based aggregators avoided divergence, whereas FedAvg and the median
diverged in \DivFedavg{} and \DivMedian{} of 120 runs.

The result changes the direction of the project. The useful channel-aware component
is erasure masking, and the reputation reward rule must be redesigned before link
telemetry can help. The findings are written up as a four-page letter in IEEE format.
The second phase will revise the reward rule, test it on fresh seeds and
higher-mobility traces, and evaluate adaptive attackers.

\textbf{Index Terms :} Federated learning, UAV swarm, Byzantine robustness,
packet erasure, hierarchical aggregation, reputation.
